% Faithful transcription of CCC2023 article15 Sections4.1–4.5. Author numbering retained explicitly.
\section*{4 A Hard Set-multilinear Polynomial in VBP}
\subsection*{4.1 Arc-partition Measure Description}
This subsection is adapted from Section 2 of [8]. Let $n=2^k$, $d\leq n$ be an even integer, and let $X=(X_1,X_2,\ldots,X_d)$ be a collection of disjoint sets of $n$ variables each. An arc-partition will be a special kind of symmetric word $w\in\{-k,k\}^d$ (i.e., a one-to-one map $\Pi$ from $X$ to $\{-k,k\}^d$). For the purpose of this subsection, the reader can even choose to think of the alphabet of $w$ as $\{-1,1\}$ (i.e., one ``positive'' and one ``negative'' value) -- we use $k,-k$ only to remain consistent with Definition 5.
Identify $X$ with the set $\{1,2,\ldots,d\}$ in the natural way. Consider the $d$-cycle graph, i.e., the graph with nodes $\{1,2,\ldots,d\}$ and edges between $i$ and $i+1$ modulo $d$. For two nodes $i\neq j$ in the $d$-cycle, denote by $[i,j]$ the arc between $i,j$, that is, the set of nodes on the path $\{i,i+1,\ldots,j-1,j\}$ from $i$ to $j$ in $d$-cycle. First, define a distribution $\mathcal D_P$ on a family of pairings (a list of disjoint pairs of nodes in the cycle) as follows. A random pairing is constructed in $d/2$ steps. At the end of step $t\in[d/2]$, we shall have a pairing $(P_1,\ldots,P_t)$ of the arc $[L_t,R_t]$. The size of $[L_t,R_t]$ is always $2t$. The first pairing contains only $P_1=\{L_1,R_1\}$ with $L_1=1$ and $R_1=2$. Given $(P_1,\ldots,P_t)$ and $[L_t,R_t]$, define the random pair $P_{t+1}$ (independently of previous choices) by
\[P_{t+1}=\begin{cases}\{L_t-2,L_t-1\}&\text{with probability }1/3\\\{L_t-1,R_t+1\}&\text{with probability }1/3\\\{R_t+1,R_t+2\}&\text{with probability }1/3.\end{cases}\]
Define $[L_{t+1},R_{t+1}]=[L_t,R_t]\cup P_{t+1}$. So, $L_{t+1}$ is either $L_t-2,L_t-1$ or $L_t$, each value is obtained with probability $1/3$, and similarly (but not independently) for $R_{t+1}$.
The final pairing is $P=(P_1,P_2,\ldots,P_{d/2})$. Denote by $P\sim\mathcal D_P$ a pairing distributed according to $\mathcal D_P$.
Once a pairing $P$ has been obtained, a word $w\in\{-k,k\}^d$ is obtained by simply randomly assigning $+k$ and $-k$ to the indices of any pair $P_i$. More formally, for every $t\in[n/2]$, if $P_t=\{i_t,j_t\}$, let with probability $1/2$, independently of all other choices,
\[w_{i_t}=+k\text{ and }w_{j_t}=-k,\]
and with probability $1/2$,
\[w_{i_t}=-k\text{ and }w_{j_t}=+k.\]
Denote by $w\sim\mathcal D$ a word in $\{-1,1\}^n$ that is sampled using this procedure. We call such a word an arc-partition. For a pair $P_t=\{i_t,j_t\}$, we refer to $i_t$ and $j_t$ as partners.
\paragraph{Definition 16 (Arc-full-rank).} We say that a polynomial $f$ that is set-multilinear over $X=(X_1,\ldots,X_d)$ is arc-full-rank if for every arc-partition $w\in\{-k,k\}^d$, $\operatorname{relrk}_w(f)=1$.
\subsection*{4.2 Construction of an Arc-full-rank Polynomial}
Below, we describe a simple construction of an ABP that computes an arc-full-rank set-multilinear polynomial. The high-level idea is to construct an ABP in which every path between start-node and end-node corresponds to a specific execution of the random process which samples arc-partitions. Each node in the ABP corresponds to an arc $[L,R]$, which sends an edge to each of the nodes $[L-2,R]$, $[L-1,R+1]$ and $[L,R+2]$. The edges have specially chosen labels that help guarantee full rank with respect to every arc-partition. For simplicity of presentation, we allow the edges of the program to be labeled by degree three polynomials in three variables. This assumption can be easily removed by replacing each edge with a constant-size ABP computing the corresponding degree three polynomial.
Formally, the nodes of the program are even-size arcs in the $d$-cycle, $d$ an even integer. The start-node of the program is the empty arc $\emptyset$ and the end-node is the whole cycle $[d]$ (both are ``special'' arcs). Let $X=(X_1,\ldots,X_d)$ be a collection of sets of variables where each $|X_i|=n$, and similarly, let $Y=(Y_1,\ldots,Y_d)$ be a distinct collection of sets of variables where each $|Y_i|=n$ (we shall refer to the $Y$-variables as auxiliary variables). For $i$ and $j$ in $\{1,\ldots,d\}$, let $X_i\cdot X_j$ denote the inner-product quadratic form $\sum_{k=1}^n x_{ik}x_{jk}$. Here, we shall assume that $X_i=\{x_{i,1},\ldots,x_{i,n}\}$ and $Y_i=\{y_{i,1},\ldots,y_{i,n}\}$.
Construct the branching program by connecting a node/arc of size $2t$ to three nodes/arcs of size $2t+2$. For $t=1$, there is just one node $[1,2]$, and the edge from start-node to it is labeled $y_{1,2}y_{2,1}(X_0\cdot X_1)$. For $t>1$, the node $[L,R]\supset[1,2]$ of size $2t<d$ is connected to the three nodes: $[L-2,R]$, $[L-1,R+1]$, and $[L,R+2]$. (It may be the case that the three nodes are the end-node.) The edge labeling is:
\begin{itemize}
\item The edge between $[L,R]$ and $[L-2,R]$ is labeled $y_{L-2,L-1}y_{L-1,L-2}(X_{L-2}\cdot X_{L-1})$.
\item The edge between $[L,R]$ and $[L-1,R+1]$ is labeled $y_{L-1,R+1}y_{R+1,L-1}(X_{L-1}\cdot X_{R+1})$.
\item The edge between $[L,R]$ and $[L,R+2]$ is labeled $y_{R+1,R+2}y_{R+2,R+1}(X_{R+1}\cdot X_{R+2})$.
\end{itemize}
Consider the ABP thus described, and the polynomial $G=G_{n,d}$ it computes. For every path $\gamma$ from start-node to end-node in the ABP, the list of edges along $\gamma$ yields a pairing $P$; every edge $e$ in $\gamma$ corresponds to a pair $P_e=\{i_e,j_e\}$ of nodes in $d$-cycle. Thus,
\[G=\sum_\gamma\prod_{e=\{i_e,j_e\}\in\gamma}y_{i_e,j_e}y_{j_e,i_e}\cdot(X_{i_e}\cdot X_{j_e}).\tag{1}\]
where the sum is over all paths $\gamma$ from start-node to end-node.
\paragraph{Remark 17.} There is in fact a one-to-one correspondence between pairings $P$ and such paths $\gamma$ (this follows by induction on $t$). Note that this is true only because pairings are tuples i.e., they are ordered by definition. Otherwise, it is of course still possible to obtain the same set of pairs in a given pairing using multiple different orderings. The sum defining $G$ can be thought of, therefore, as over pairings $P$.
The following statement summarizes the main useful property of $G$.
\paragraph{Lemma 18.} Over any field $\mathbb F$ of characteristic zero, the polynomial $G=G_{n,d}$ defined above is arc-full-rank as a set-multilinear polynomial in the variables $X$ over the field $\mathbb F(Y)$ of rational functions in $Y$.
\paragraph{Proof.} Let $w\sim\mathcal D$ be an arc-partition. We want to show that $M_w(G)$ has full rank. The arc-partition $w$ is defined from a pairing $P=(P_1,\ldots,P_{d/2})$ (though as discussed in Remark 17, there could be multiple such $P$). The pairing $P$ corresponds to a path $\gamma$ from start-node to end-node. Consider the polynomial $f$ that is obtained by setting every $y_{i,j}=y_{j,i}=0$ in $F$ such that $\{i,j\}$ is not a pair in $P$, and setting every $y_{i,j}=y_{j,i}=1$ for every pair $\{i,j\}$ in $P$. Then, it is easy to see that the only terms that survive in (1) correspond to paths (and in turn, pairings) which have the same underlying set of pairs as $P$. As a consequence, $f$ is simply some non-zero constant times a polynomial which is full-rank. $M_w(f)$ being full rank then implies that $M_w(G)$ is also full-rank\footnote{This argument is the same as in the proof of Lemma 15.}.\hfill$\triangleleft$
\subsection*{4.3 Bounding $\operatorname{relrk}_w$ for Small Set-multilinear Formulas}
As discussed in Section 1, the high-level strategy to prove Theorem 1 is to show that the relative rank (with respect to arc-partitions) of our hard polynomial is large (as already established in Lemma 18), while it is small for (small enough) set-multilinear formulas. The remainder of the section is devoted to establishing the latter. Before moving on to it, we shall first state the following more precise and general version of Theorem 1. We also incorporate Remark 2 here and show our lower bound for any degree $d\leq n$. Theorem 1 follows from the special case $d=n$.
\paragraph{Theorem 19.} Let $n=2^k$, and suppose $d\leq n$ be an even integer that is large enough\footnote{We only need $d$ to be larger than some absolute constant.}, and $1\leq\Delta\leq\log d/\log\log d$ be any integer. Let $X_i,Y_i$ denote the sets of $n$ variables $\{x_{i,j}:j\in[n]\}$ and $\{y_{i,j}:j\in[n]\}$ respectively and let $X,Y$ be the tuples $(X_1,\ldots,X_d)$ and $(Y_1,\ldots,Y_d)$. Then,
\begin{itemize}
\item there is a $\operatorname{poly}(n,d)$-size branching program computing $G_{n,d}$ as defined above whose every internal node computes a set-multilinear polynomial,
\item any set-multilinear formula of product-depth $\Delta$ computing $G_{n,d}$ must have size at least $d^{\Omega(d^{1/\Delta}/\Delta)}$.
\item Further, any set-multilinear formula of arbitrary product-depth computing $G_{n,d}$ must have size at least $d^{\Omega(\log d)}$.
\end{itemize}
The following couple of lemmas formalize the high-level idea mentioned before the statement of Theorem 19 -- they correspond to the low-depth case and general depth case respectively. Most of the remainder of this section is devoted to the proof of Lemma 20; Lemma 22 has a similar (and in fact, easier) proof and for this reason, we only provide a sketch that is deferred to the appendix.
\paragraph{Lemma 20.} Let $\mathbb K$ be any field and let $X_1,\ldots,X_D$ be sets of $n$ distinct variables each. Let $C$ be a set-multilinear formula over $\mathbb K$ of constant product-depth $\Delta\geq1$ of size at most $s$ which computes a polynomial over $\mathbb K$ that is set-multilinear with respect to the partition $(X_{i_1},\ldots,X_{i_d})$ where $1\leq i_1<\cdots<i_d\leq D$ and each $|X_i|=n$. Let $w\sim\mathcal D$ be an arc-partition sampled from $\{-k,k\}^D$. Then, we have
\[\operatorname{relrk}_w(C)\leq s\cdot2^{-kd^{1/\Delta}/2000}\]
with probability at least $1-s\cdot d^{-d^{1/\Delta}/(10^7\Delta)}$.
\paragraph{Remark 21.} Note that in the statement above, we are abusing notation and overloading the $\operatorname{relrk}_w$ notation -- assume that $\operatorname{relrk}_w(C)$ is defined in the obvious projective manner i.e., if $S=\{i_1,\ldots,i_d\}$, then $\operatorname{relrk}_w(C):=\operatorname{relrk}_{w|_S}(C)$ where $w|_S$ is as defined in Definition 6.
\paragraph{Lemma 22.} Let $\mathbb K$ be any field and let $X_1,\ldots,X_d$ be sets of $n$ distinct variables each. Let $F$ be a set-multilinear formula over $\mathbb K$ of size at most $s$ which computes a polynomial over $\mathbb K$ that is set-multilinear with respect to the partition $(X_1,\ldots,X_d)$ where each $|X_i|=n$. Let $w\sim\mathcal D$ be an arc-partition sampled from $\{-k,k\}^d$. Then, we have
\[\operatorname{relrk}_w(F)\leq s\cdot2^{-k\log d/2000}\]
with probability at least $1-s\cdot d^{-\log d/(10^7\Delta)}$.
Before moving on to the technical core of this section (the proof of Lemma 20), let us finish the proof of Theorem 19.
\paragraph{Proof of Theorem 19 given Lemmas 20 and 22.} Note that the first item follows immediately from the definition of $G_{n,d}$ (see (1)). Let us prove the last two items of Theorem 19. Let $C$ be a set-multilinear formula of product depth $\Delta$ of size $s$ computing $G_{n,d}(X)$ (now interpreted as a formula over the field $\mathbb F(Y)$). Suppose $s<d^{d^{1/\Delta}/(2\times10^7\Delta)}$. Then, by Lemma 20, for an arc-partition $w\sim\mathcal D$ sampled from $\{-k,k\}^d$, it follows that with probability at least $1-d^{-d^{1/\Delta}/(2\times10^7\Delta)}$,
\[\operatorname{relrk}_w(C)\leq s\cdot2^{-kd^{1/\Delta}/2000}.\]
Fix such an arc-partition $w$. By Lemma 18, we have
\[s\geq2^{kd^{1/\Delta}/2000}\cdot\operatorname{relrk}_w(C)=n^{d^{1/\Delta}/2000}\]
which contradicts the assumption that $s<d^{d^{1/\Delta}/(2\times10^7\Delta)}$. Thus, we conclude that $s\geq d^{d^{1/\Delta}/(2\times10^7\Delta)}=d^{\Omega(d^{1/\Delta}/\Delta)}$.
Similarly, to see the final item of Theorem 19, let $F$ be a set-multilinear formula of size $s$ computing $G_{n,d}$ (now interpreted as a formula over the field $\mathbb F(Y)$). Suppose $s<d^{\log d/(2\times10^7)}$. Then, for an arc-partition $w\sim\mathcal D$ sampled from $\{-k,k\}^d$, by Lemma 22, with probability at least $1-d^{-\log d/(2\times10^7)}$,
\[\operatorname{relrk}_w(F)\leq s\cdot2^{-k\log d/2000}.\]
Fix such an arc-partition $w$. Therefore, again by Lemma 18,
\[s\geq2^{k\log d/2000}\cdot\operatorname{relrk}_w(F)=n^{\log d/2000}\]
which contradicts the assumption that $s<d^{\log d/(2\times10^7)}$. Thus, we conclude that $s\geq d^{\log d/(2\times10^7)}=d^{\Omega(\log d)}$.\hfill$\triangleleft$
The essential ingredient in the proof of Lemma 20 is a combinatorial proposition which we will call the ``many violations lemma''. As alluded to in Section 1.5, this is a modification of a corresponding statement in [8] (Lemma 4.1). However, because we are working in the low-depth setting (as opposed to [8]) and because we are seeking such strong and near-optimal lower bounds, we need to make significant changes -- this includes introducing new conceptual arguments to tighten the analysis. To state this lemma, we shall reproduce some of the definitions made in Section 4 of [8].
Again, we identify the set of variables $X=(X_1,\ldots,X_D)$ with the $D$-cycle $\{1,\ldots,D\}$, where addition is modulo $D$. Let $S$ be a collection of disjoint subsets of the cycle to $K$ parts, namely, $S=(S_1,\ldots,S_K)$ where each $S_k\subset\{1,\ldots,D\}$ and $S_k\cap S_{k'}=\emptyset$ for all $k\neq k'$ in $[K]$. We also think of $[K]$ as a set of colors, and of $S$ as a (partial) $K$-coloring of some $d$ nodes of the cycle, where $d=|S_1|+\cdots+|S_K|$. We shall refer to the nodes in the $D$-cycle outside of $S$ as uncolored.
For a pairing $P$, define the number of $k$-violations by
\[V_k(P)=\{P_t\in P:|P_t\cap S_k|=1\}.\]
In words, it is the set of pairs in which one color is $k$ and the other color is different. Fix $\varepsilon=1/1000$ and denote
\[G(P)=\{k\in[K]:|V_k(P)|\geq d^\varepsilon\}.\]
We do not include $S$ as a subscript in these two notations since $S$ will be known from the context (and will be fixed throughout most of the discussion). The next crucial lemma shows that for every fixed non-redundant $K$-coloring of the cycle, a random pairing has, with high probability, many colors with many violations.
\paragraph{Lemma 23 (Many Violations Lemma).} For all large enough $d$ and for all integers $K$ in the range $[2d^{1/(\Delta+1)}/3,2d^{1/(\Delta+1)}]$ the following holds: Let $S=(S_1,\ldots,S_K)$ be a collection of disjoint subsets of the $D$-cycle and suppose that $|S_k|\geq d^{\Delta/(\Delta+1)}/2$ for all $k\in[K]$. Then,
\[\mathbb P[G(P)\leq K/1000]\leq d^{-K/500},\]
where $P\sim\mathcal D_P$.
\paragraph{Remark 24.} Other than the differences in parameter ranges, one key difference between the statement above from Lemma 4.1 in [8] is the loosening of the requirement that $S$ be a partition of the $D$-cycle. Note that here, we only demand that $S$ be a collection of disjoint subsets (i.e., some nodes are allowed to remain uncolored) -- this requirement is indeed key for the inductive proof of Lemma 20 to go through.
Before proving Lemma 23, let us next see that the many violations lemma suffices to prove the relative rank upper bound on low-depth set-multilinear formulas.
\paragraph{Proof of Lemma 20 given Lemma 23.} We prove the statement by induction on $\Delta$. Identify the set $\{i_1,\ldots,i_d\}$ with $[d]$.
If $\Delta=1$, then $C=C_1+\cdots+C_t$ where each $C_i$ is a product of linear forms. So, for all $i\in[t]$, by Claim 7,
\[\operatorname{relrk}_w(C_i)=\prod_{i=1}^d2^{-\frac12|w_j|}=2^{-kd/2}\]
where in the last step, we used the observation that regardless of the choice of $w$, $|w_j|=k$ for all $j\in[n]$. Hence, by the sub-additivity of $\operatorname{relrk}_w$, with probability $1$, we have
\[\operatorname{relrk}_w(C)\leq s\cdot2^{-kd/2}\leq s\cdot2^{-kd/2000}.\]
Next, we assume the statement is true for all formulas of product-depth $\leq\Delta$. Let $C$ be a formula of product-depth $\Delta+1$. So, $C$ is of the form $C=C_1+\cdots+C_t$. Using a similar terminology to that in [23] and [22], we say that a sub-formula $C_i$ of size $s_i$ is of type 1 if one of its factors has degree at least $T_\Delta=d^{\Delta/(\Delta+1)}$, otherwise we say it is of type 2.
Suppose $C_i=C_{i,1}\cdots C_{i,t_i}$ is of type 1 with, say, $C_{i,1}$ having degree at least $T_\Delta$. Let $w_{i,1}$ be the corresponding word i.e., $w_{i,1}=w|_{S_1}$ if $C_{i,1}$ is set-multilinear with respect to $S_1\subsetneq[d]$. If it has size $s_{i,1}$, then since it has product-depth at most $\Delta$, it follows by induction that
\[\operatorname{relrk}_w(C_i)\leq\operatorname{relrk}_{w_{i,1}}(C_{i,1})\leq s_{i,1}\cdot2^{-kT_\Delta^{1/\Delta}/2000}\leq s_i\cdot2^{-kd^{1/(\Delta+1)}/2000}\]
with probability at least
\[1-s_{i,1}\cdot T_\Delta^{-T_\Delta^{1/\Delta}/(10^7\Delta)}\geq1-s_i\cdot d^{-\frac{d^{1/(\Delta+1)}}{10^7\Delta}\cdot\frac{\Delta}{\Delta+1}}=1-s_i\cdot d^{-d^{1/(\Delta+1)}/(10^7(\Delta+1))}.\]
Now suppose that $C_i=C_{i,1}\cdots C_{i,t_i}$ is of type 2 i.e., each factor $C_{i,j}$ has degree $<T_\Delta$. Note that this forces $t_i>d/T_\Delta=d^{1/(\Delta+1)}$. As the formula is set-multilinear, $(S_1,\ldots,S_{t_i})$ form a partition of $[d]$ where each $C_{i,j}$ is set-multilinear with respect to $(X_\ell)_{\ell\in S_j}$ and $C_i$ is set-multilinear with respect to $(X_\ell)_{\ell\in S}$. Let $w_{i,1},\ldots,w_{i,t_i}$ be the corresponding decomposition, whose respective sums are denoted simply by $w_{S_1},\ldots,w_{S_{t_i}}$.
From the properties of $\operatorname{relrk}_w$ (Claim 7), we have
\[\operatorname{relrk}_w(C_i)=\prod_{j=1}^{t_i}\operatorname{relrk}_{w_{i,j}}(C_{i,j})\leq\prod_{j=1}^{t_i}2^{-\frac12|w_{S_j}|}=2^{-\frac12\sum_{j=1}^{t_i}|w_{S_j}|},\]
from which we observe that the task of upper bounding $\operatorname{relrk}_w(C)$ can be reduced to the task of lower bounding the sum $\sum_{j=1}^{t_i}|w_{S_j}|$, which is established in the following claim. For the sake of convenience, the choice of the alphabet for $w$ below is scaled down to $\{-1,1\}$.
\paragraph{Claim 25.} For large enough $d$, suppose $(S_1,\ldots,S_K)$ is a partition of $[d]$ such that each $|S_j|<T_\Delta=d^{\Delta/(\Delta+1)}$. Then, we have
\[\mathbb P_{w\sim\mathcal D}\left[\sum_{j=1}^K|w_{S_j}|<\frac{d^{1/(\Delta+1)}}{2000}\right]\leq d^{-d^{1/(\Delta+1)}/10^7}.\]
Here, $\mathcal D$ refers to the original distribution i.e., an arc-partition over the $D$-cycle.
\paragraph{Proof.} We first show that without loss of generality, we may assume that each $S_j$ has size ``roughly'' $T_\Delta$. To see this, we apply the following clubbing procedure to the sets in the partition $(S_1,\ldots,S_K)$:
\begin{itemize}
\item Start with the given partition $(S_1,\ldots,S_K)$. At each step in the procedure, we shall ``club'' two of the sets in the partition according to the following rule.
\item If there are two distinct sets $S'$ and $S''$ in the current partition each of size $<T_\Delta/2$, we remove both of them and add their union $S'\cup S''$ to the partition.
\item If the rule above is no longer applicable, then we have at most one set in the current partition of size $<T_\Delta/2$. If there is none, then we halt the procedure here. Otherwise, we union this set with any one of the other sets and then halt.
\end{itemize}
After the clubbing procedure, we are left with a partition $(S'_1,\ldots,S'_{K'})$ of $[d]$ such that $T_\Delta/2\leq|S'_j|\leq3T_\Delta/2$ for each $j\in[K']$, also implying that $2d^{1/(\Delta+1)}/3\leq K'\leq2d^{1/(\Delta+1)}$. Through a repeated use of the triangle inequality, we see that $\sum_{j=1}^{K'}|w_{S'_j}|\leq\sum_{j=1}^K|w_{S_j}|$. Therefore, upper bounding the latter sum is a ``smaller'' event than upper bounding the former sum. Hence, it suffices to prove the statement of the claim with the assumption that $T_\Delta/2\leq|S_j|\leq3T_\Delta/2$ for each $j\in[K]$ (we henceforth drop the primed notation).
Applying Lemma 23 to the tuple $(S_1,\ldots,S_K)$, we obtain that
\[\mathbb P[G(P)\leq K/1000]\leq d^{-K/500}.\]
The idea is to condition on the high probability event that $G(P)>K/1000$. Fix a pairing $P$ with this property. Consider an ordering $\sigma$ of the colors in $G(P)$. A color $k$ is said to be bright with respect to an ordering if there are at least $d^\varepsilon/2$ nodes $x$ of color $k$ such that either the partner of $x$ is uncolored or its partner is colored using a color that appears after $k$ in the ordering $\sigma$. Call an ordering $\sigma$ of the nodes in $G(P)$ good if there are at least $|G(P)|/2$ bright colors with respect to $\sigma$. The observation is that for any ordering $\sigma$ of the colors, either $\sigma$ itself is good, or its reverse is good. We conclude that given any pairing $P$, there exists a good ordering of $G(P)$. Fix any such good ordering and let $H(P)$ be the collection of bright colors with respect to this ordering.
Next, notice that if the sum $\sum_{j=1}^K|w_{S_j}|$ is at most $d^{1/(\Delta+1)}/2000$, then so is the sum $\sum_{k\in H(P)}|w_{S_k}|$. Let $K'=|H(P)|$ (which is at least $K/2000$ if $G(P)>K/1000$). View the sampling of $\Pi$ from $P$ as happening in a specific order, according to the order of $k_1,k_2,\ldots,k_{K'}$: First define $\Pi$ on pairs with at least one point with color $k_1$, then define $\Pi$ on remaining pairs with at least one point with color $k_2$, and so forth. When finished with $k_1,\ldots,k_{K'}$, continue to define $\Pi$ on all other pairs.
Conditioned on the event that $G(P)>K/1000$, this implies that $|w_{S_j}|\leq1$ for each $j\in H(P)$. For every $j\in H(P)$, define $E_j$ to be the event that $|w_{S_{k_j}}|\leq1$. By choice, conditioned on $E_1,\ldots,E_{j-1}$, there are at least $d^\varepsilon/2$ pairs $P_t$ so that $|P_t\cap S_{k_j}|=1$ that are not yet assigned a ``positive'' or ``negative'' sign. For every such $P_t$, the element in $P_t\cap S_{k_j}$ is assigned a positive sign with probability $1/2$, and is independent of any other $P_{t'}$. The probability that a binomial random variable $B$ over a universe of size $U\geq d^\varepsilon/2$ and marginals $1/2$ obtains any specific value is at most $O(U^{-1/2})=O(d^{-\varepsilon/2})$. Hence, for all $j\in H(P)$, by the union bound,
\[\mathbb P[E_j\mid E_1,\ldots,E_{j-1},P]\leq\mathbb P_B[U/2-1\leq B\leq U/2+1]\leq O(3\cdot d^{-\varepsilon/2})\leq d^{-\varepsilon/4}.\]
Therefore,
\[\mathbb P[|w_{S_{k_j}}|\leq1\text{ for all }j\in H(P)]\leq\mathbb E[d^{-\varepsilon|H(P)|/4}\mid G(P)>K/1000]+d^{-K/500}\leq d^{-K/10^7}.\]
Finally, we note that
\[\mathbb P_{w\sim\mathcal D}\left[\sum_{j=1}^K|w_{S_j}|<\frac{d^{1/(\Delta+1)}}{2000}\right]\leq\mathbb P[|w_{S_{k_j}}|\leq1\text{ for all }j\in H(P)].\]\hfill$\triangleleft$
The claim above and the preceding calculation immediately implies that for a sub-formula $C_i$ of type 2,
\[\operatorname{relrk}_w(C_i)\leq s_i\cdot2^{-kd^{1/(\Delta+1)}/2000}\]
with probability at least $1-d^{-d^{1/(\Delta+1)}/10^7}\geq1-s_i\cdot d^{-d^{1/(\Delta+1)}/(10^7(\Delta+1))}$.
Next, by a union bound over $i\in[t]$ and the sub-additivity property of $\operatorname{relrk}_w$, it follows that
\[\operatorname{relrk}_w(C)\leq\operatorname{relrk}_w(C_1)+\cdots+\operatorname{relrk}_w(C_t)\leq s_1\cdot2^{-kd^{1/(\Delta+1)}/2000}+\cdots+s_t\cdot2^{-kd^{1/(\Delta+1)}/2000}=s\cdot2^{-kd^{1/(\Delta+1)}/2000}\]
with probability at least $1-s\cdot d^{-d^{1/(\Delta+1)}/(10^7(\Delta+1))}$, which concludes the proof of the lemma.\hfill$\triangleleft$

\subsection*{4.4 Proof of the Many Violations Lemma}
Fix some collection of disjoint subsets (or a ``partial'' coloring) $S=(S_1,\ldots,S_K)$ of the $D$-cycle satisfying the conditions of the lemma. Think of $S$ as a partial function from the $D$-cycle to the set $[K]$, either assigning a node its color in $[K]$ or leaving it uncolored; $S(i)$ is the color of $i$. Use the following definition to partition the proof into cases. For a color $k$, count the number of jumps in it (with respect to the partition $S$) to be
\[J_k=\{j\in S_k:k=S(j)\neq S(j+1)\},\]
the set of elements $j$ of color $k$ so that $j+1$ is either uncolored or has a color different from $k$. As mentioned previously, this subsection is adapted from the proof of Lemma 4.1 in [8]. In what follows, we include remarks where we require a more refined analysis than [8] or a different argument to suit the parameter demands of Lemma 23. Overall, we have attempted to provide a more comprehensive and complete exposition to the proof of the many violations lemma.
\paragraph{Case 1: Many colors with many jumps.}
The high-level idea is that each color with many jumps has many violations because pairs of the form $(j,j+1)$ yield violations as soon as they are constructed.
Assume that for at least $K/2$ colors $k$, $|J_k|>d^{2\varepsilon}$. Denote by $B\subseteq[K]$ the set of $k$'s that satisfy this inequality. Then, for every $k$ in $B$, there exists a subset $Q_k\subset J_k$ of size $N=\lceil d^{2\varepsilon}\rceil$. Let
\[Q:=\bigcup_{k\in B}Q_k.\]
We think of the construction of the (random) pairing $P$ as happening in epochs, depending on $Q$, as follows.
For $t>0$, define the random variable
\[Q(t)=Q\setminus[L_t-4,R_t+4],\]
the set $Q$ after removing a four-neighborhood of $[L_t,R_t]$. For a certain sequence of time steps $t$, we will define special nodes $q_t$ which lie in this small ``cloud'' around the arc $[L_t,R_t]$ (i.e., within a distance of 4 on either side of the arc) -- it is for these special nodes $q_t$ that the set of pairs $(q_t,q_{t+1})$ will provide many violations. We now formalize this intuition.
Let $\tau_1\geq\tau_0:=1$ be the first time $t$ after $\tau_0$ so that the distance between $[L_t,R_t]$ and $Q(\tau_0)$ is at most two. The distance between $[L_{\tau_0},R_{\tau_0}]$ and $Q(\tau_0)$ is at least five. The size of the arc $[L_t,R_t]$ increases by two at each time step. So, $\tau_1\geq\tau_0+2$. Let $q_1$ be an element of $Q(\tau_0)$ that is of distance at most two from $[L_{\tau_1},R_{\tau_1}]$; if there is more than one such $q_1$, choose arbitrarily. The minimality of $\tau_1$ implies that $q_1$ is not in $[L_{\tau_1},R_{\tau_1}]$.
Let $\tau_2\geq\tau_1$ be the first time $t$ after $\tau_1$ so that the distance between $[L_t,R_t]$ and $Q(\tau_1)$ is at most two. Let $q_2$ be an element of $Q(\tau_1)$ that is of distance at most two from $[L_{\tau_2},R_{\tau_2}]$. Define $\tau_j,q_j$ for $j>2$ similarly, until $Q(\tau_j)$ is empty. As long as $|Q(\tau_j)|\geq8$, we have $|Q(\tau_j+1)|\geq|Q(\tau_j)|-8$. This process, therefore, has at least $KN/16$ steps. For $1\leq j\leq KN/16$, denote by $E_j$ the event that during the time between $\tau_j$ and $\tau_{j+1}$ the pair $\{q_j,q_j+1\}$ is added to $P$. The pair $\{q_j,q_j+1\}$ is violating color $S(q_j)$. At time $\tau_j$, even conditioned on all the past $P_1,\ldots,P_{\tau_j}$, in at most two steps (and before $\tau_{j+1}$) we can add the pair $\{q_j,q_j+1\}$ to $P$. For every $j$, therefore,
\[\mathbb P[E_j\mid P_1,\ldots,P_{\tau_j}]\geq(1/3)(1/3)=1/9.\]
Next, let $N'=\lceil KN/960\rceil$. We want to show that with high probability, for at least $N'$ many $j$, the event $E_j$ occurs. There are $\binom{\lfloor KN/16\rfloor}{\lceil KN/960\rceil}$ many ways of choosing a set of indices $j$ of size $N-N'$. Subsequently,
\begin{align*}
\mathbb P[\text{there is }j_1,\ldots,j_{N'}\text{ so that }E_{j_1}\cap\cdots\cap E_{j_{N'}}]
&\geq1-\binom{\lfloor KN/16\rfloor}{\lceil KN/960\rceil}\cdot(8/9)^{N-N'}\\
&\geq1-(960e/16)^{N'}\cdot(8/9)^{60N'}\\
&\geq1-c^{N'}
\end{align*}
where $0<c<1$ is a universal constant. Finally, we argue that if there do exist $j_1,\ldots,j_{N'}$ for which the events $E_{j_1},\ldots,E_{j_{N'}}$ occur, then $G(P)\geq K/1000$. To see this, note that the size of every $Q_k$ is $N$. So, every color $k$ in $B$ can contribute at most $N$ elements to $j_1,\ldots,j_{N'}$. If $G(P)<K/1000$, then at most these many colors can contribute larger than $d^\varepsilon$ (and up to $N$ elements) -- combined, at most $KN/1000$ elements. However, there are at least $K/2-K/1000$ colors which can contribute only up to $d^\varepsilon$ elements. Again combined, this is not sufficient to cover the $N'$ elements overall (for large enough $d$), which is a contradiction. Hence,
\[\mathbb P[G(P)\geq K/1000]\geq\mathbb P[\text{there is }j_1,\ldots,j_{N'}\text{ so that }E_{j_1}\cap\cdots\cap E_{j_{N'}}].\]
and the proof follows in this case as $c^{N'}\ll d^{-\Omega(K)}$.
\paragraph{Case 2: Many colors with few jumps.}
The intuition is that many violations will come from pairs of the form $\{L_t-1,R_t+1\}$ in the construction of the pairing. Assume that for at least $K/2$ colors $k$, $|J_k|\leq d^{2\varepsilon}$. Denote again by $B\subseteq[K]$ the set of $k$'s that satisfy the above inequality. We say that a color $k$ is noticeable in the arc $A$ if
\[d^{\Delta/(\Delta+1)-4\varepsilon}\leq|S_k\cap A|\leq|A|-d^{\Delta/(\Delta+1)-4\varepsilon}.\]
\paragraph{Claim 26.} There are $K'\geq K/2-1$ disjoint arcs $A_1,\ldots,A_{K'}$ so that for every $j\in[K']$,
\begin{enumerate}
\item $|A_j|=m=\lfloor d^{\Delta/(\Delta+1)-3\varepsilon}\rfloor$ and,
\item there is a color $k_j$ in $B$ that is noticeable in $A_j$.
\end{enumerate}
Moreover, the colors $k_1,\ldots,k_{K'}$ can be chosen to be pairwise distinct.
\paragraph{Proof.} For each color $k$ in $B$, there are at least $d^{\Delta/(\Delta+1)}/2$ vertices of color $k$ in the $D$-cycle and at most $d^{2\varepsilon}$ jumps in the color $k$. Therefore, there is at least one $k$-monochromatic arc of size at least $d^{\Delta/(\Delta+1)-2\varepsilon}$. Hence, on the $D$-cycle, there are such monochromatic arcs $I_{k_1},\ldots,I_{k_{|B|}}$ for the colors $k_1,\ldots,k_{|B|}$ in $B$, in this order $(1<2<\cdots<D)$.
Consider an arc $A$ of size $m$ included in $I_{k_1}$. Thus $|S_{k_1}\cap A|=m$. If we ``slide'' the arc $A$ until it is included in $I_{k_2}$, then $|S_{k_1}\cap A|=0$. By continuity, there is an intermediate position for the arc $A$ such that $d^{\Delta/(\Delta+1)-4\varepsilon}\leq|S_{k_1}\cap A|\leq m-d^{\Delta/(\Delta+1)-4\varepsilon}$. This provides the first arc $A_1$ of the claim.
Sliding an arc inside $I_{k_2}$ to inside $I_{k_3}$ shows that there exists an arc $A_2$ such that $d^{\Delta/(\Delta+1)-4\varepsilon}\leq|S_{k_2}\cap A_2|\leq m-d^{\Delta/(\Delta+1)-4\varepsilon}$. The arcs $A_1$ and $A_2$ are disjoint: The distance of the largest element of $A_1$ and the smallest element of $I_{k_2}$ is at most $m$. The distance of the smallest element of $A_2$ and the largest element of $I_{k_2}$ is at most $m$. The size of $I_{k_2}$ is larger than $2m$. Proceed in this way to define $A_3,\ldots,A_{|B|-1}$.\hfill$\triangleleft$
Use Claim 26 to divide the construction of the (random) pairing into epochs. Denote by $A^{(0)}$ the family of arcs given by the claim. Let $\tau_1$ be the first time $t$ that the arc $[L_t,R_t]$ hits one of the arcs in $A^{(0)}$. Denote by $A_1$ that arc that is hit at time $\tau_1$ (break ties arbitrarily). Denote by $k_1$ the color that is noticeable in $A_1$. Let $\sigma_1$ be the first time $t$ so that $A_1$ is contained in $[L_t,R_t]$. Let $A^{(1)}$ be the subset of $A^{(0)}$ of arcs that have an empty intersection with $[L_{\sigma_1},R_{\sigma_1}]$. Similarly, let $\tau_2$ be the first time $t$ after $\sigma_1$ that the arc $[L_t,R_t]$ hits one of the arcs in $A^{(1)}$. If there are no arc in $A^{(1)}$, define $\tau_2=\infty$. Denote by $A_2$ that arc that is hit at time $\tau_2$. Denote by $k_2$ the color that is noticeable in $A_2$. Let $\sigma_2$ be the first time $t$ so that $A_2$ is contained in $[L_t,R_t]$. Let $A^{(2)}$ be the subset of $A^{(1)}$ of arcs that have an empty intersection with $[L_{\sigma_2},R_{\sigma_2}]$. Define $\tau_j,\sigma_j,A_j,k_j,A^{(j)}$ for $j>2$ analogously. For every $j\geq1$, denote by $E_j$ the event that during the time between $\tau_j$ and $\tau_{j+1}$ the number of pairs added that violate color $k_j$'s at most $d^\varepsilon$. (If $E_j$ does not hold, then $|V_{k_j}(P)|\geq d^\varepsilon$ and $k_j\in G(P)$. The main part of the proof is summarized in the following proposition, whose proof is deferred to Section 4.5.
\paragraph{Lemma 27 (Chessboard Lemma).} Let $\delta=0.10$. For every $j\geq1$, and any choice of pairs $P_1,\ldots,P_{\tau_j}$,
\[\mathbb P[E_j\mid P_1,\ldots,P_{\tau_j},|A^{(j-1)}|\geq3]\leq d^{-\delta}.\]
Given this lemma, let us finish the proof of Lemma 23. Define $K''=\lfloor K'/10\rfloor$ and let $T$ denote the event that the number of $j$'s for which $|A^{(j)}|\geq3$ is at least $K''$. First, we argue that $T$ occurs with high probability.
For any $j\geq1$, consider the evolution of the arc $[L_t,R_t]$ between the time steps $\tau_j$ (when it first hits arc $A_j$) and $\sigma_j$ (when it completely engulfs it). During this epoch, let us call the evolution of $[L_t,R_t]$ in the ``direction'' of $A_j$ as good (labelled ``G'') and away from the direction of $A_j$ as bad (``B''). To this end, for any time step in this epoch, we can code the three possible choices for the evolution of $[L_t,R_t]$ as GG (when the arc is grown in the direction of $A_j$), GB (when the arc is grown equally on either side), or BB (when it is grown away from the direction of $A_j$). Consequently, the evolution of $[L_t,R_t]$ during this epoch can be realized as a sequence consisting of the symbols G and B.
Consider the sequence $s$ of G's and B's obtained by concatenating the sequences corresponding to all the epochs (ignoring the choices made at time steps that do not lie in such epochs, i.e., between $\tau_j$ and $\sigma_j$ for some $j$ -- as there is no corresponding notion of a ``good'' direction outside such epochs). The intuition is that if $|A^{(K'')}|<3$ (i.e., if $T$ does not occur), then there must be an extremely large number of B's compared to G's (i.e., the arc $[L_t,R_t]$ evolves disproportionately in the bad direction) in the concatenated string $s$, which should occur only with a vanishingly small probability.
Consider the sub-string $s'$ of $s$ that corresponds to the choices made only for the nodes in $A^{(0)}\setminus A^{(K'')}$. Note that there are precisely $mK''$ many G's in $s'$. Suppose $|A^{(K'')}|=2$ for concreteness (the cases $|A^{(K'')}|=1$ and $|A^{(K'')}|=0$ are similar). This implies that there are $m(K'-2-K'')$ many B's in $s'$. Since only up to $mK''$ many of these B's may appear as a result of the evolution making a choice of the form GB, it follows that the evolution of $[L_t,R_t]$ must make a choice of the form BB at least $m(K'-2-2K'')/2$ times out of a possible $m(K'-2)/2$, in order to cover the elements of $A^{(0)}\setminus A^{(K'')}$. Denote $K_1:=(K'-2)/2$. By the union bound, this probability is at most
\[\mathbb P[|A^{(K'')}|=2]\leq\binom{mK_1}{mK''}\cdot(1/3)^{m(K_1-K'')}<c_2^{mK''}\]
for some universal constant $0<c_2<1$. Similarly, we have bounds for both $\mathbb P[|A^{(K'')}|=1]$ and $\mathbb P[|A^{(K'')}|=0]$ and it follows that $\mathbb P[T]\geq1-c^{mK''}$ for some universal constant $0<c<1$.
\paragraph{Remark 28.} The argument above for showing that $T$ occurs with high probability differs considerably from [8], where the corresponding event is sketched to occur with probability only at least $1-d^{cm^{1/3}}$, which is not strong enough for our purposes.
Next, note that
\[\mathbb P[G(P)<K/1000]\leq\mathbb P[G(P)<K/1000\cap T]+\mathbb P[\neg T]\leq\mathbb P[G(P)<K/1000\mid T]+\mathbb P[\neg T].\]
If $G(P)<K/1000$, then at least $K/2-K/1000$ colors in $B$ have at most $d^\varepsilon$ many violations. Since $K''=\lfloor K'/10\rfloor<K/2-K/1000$, in particular, there must exist at least $K''/2$ colors within the first $K''$ colors (here we are using the ordering of colors as provided by Claim 26) for which there are at most $d^\varepsilon$ many violations. We then obtain the following by conditioning on $T$, using the union bound.
\[\mathbb P[G(P)<K/1000\cap T]\leq2^{K''}\max_{H=\{j_1<\cdots<j_{K''/2}\}\subset[K'']}\mathbb P[E_{j_1},\ldots,E_{j_{K''/2}}\mid|A^{(K'')}|\geq3].\]
For a fixed choice of $H$, by the chain rule and Lemma 27, we have
\begin{align*}
&\mathbb P[E_{j_1}\cap\cdots\cap E_{j_{K''/2}}\mid|A^{(K'')}|\geq3]\\
&=\mathbb P[E_{j_1}\mid T]\cdot\mathbb P[E_{j_2}\mid E_{j_1}\cap T]\cdots\mathbb P[E_{j_{K''/2}}\mid E_{j_{K''/2-1}}\cap\cdots\cap E_{j_1}\cap T]\\
&\leq d^{-\delta K''/2}\leq d^{-0.1K'/20}\leq d^{-K/400}.
\end{align*}
Overall, we conclude that
\[\mathbb P[G(P)<K/1000]\leq d^{-K/500}.\]
\subsection*{4.5 Proof of the Chessboard Lemma}
To prove Lemma 27, we use a different point of view of the random process. We begin by describing this different view, and later describe its formal connection to the distribution on pairings. This subsection is adapted from Section 5 of [8] and closely follows their argument, though with numerous parameter changes to suit our demands.
The view uses two definitions. One is a standard definition of a two-dimensional random walk, and the other is a definition of a ``chessboard'' configuration in the plane. The proof of the proposition will follow by analyzing the behavior of the random walk on the ``chessboard''. Let $d$ be as above and $m$ be as defined in Lemma 26. The random walk $W$ on $\mathbb N^2$ is defined as follows. It starts at the origin, $W_0=(0,0)$. At every step it move to one of three nodes, independently of previous choices,
\[W_{t+1}=\begin{cases}W_t+(0,2)&\text{with probability }1/3\\W_t+(1,1)&\text{with probability }1/3\\W_t+(2,0)&\text{with probability }1/3.\end{cases}\]
At time $t$, the $L_1$-distance of $W_t$ from the origin is thus $2t$.
The ``chessboard'' is defined as follows. Let $\alpha_1:[m]\to\{0,1\}$ and $\alpha_2:[2m]\to\{0,1\}$ be two Boolean functions. The functions $\alpha_1,\alpha_2$ induce a ``chessboard'' structure on the board $[m]\times[2m]$. A position in the board $\xi=(\xi_1,\xi_2)$ is colored either white or black. It is colored black if $\alpha_1(\xi_1)\neq\alpha_2(\xi_2)$ and white if $\alpha_1(\xi_1)=\alpha_2(\xi_2)$. We say that the ``chessboard'' is well-behaved if
\begin{enumerate}
\item $\alpha_1$ is far from constant:
\[d^{\Delta/(\Delta+1)-4\varepsilon}\leq|\{\xi_1\in[m]:\alpha_1(\xi_1)=1\}|\leq m-d^{\Delta/(\Delta+1)-4\varepsilon}.\]
\item $\alpha_1$ does not contain many jumps:
\[|\{\xi_1\in[m-1]:\alpha_1(\xi_1)\neq\alpha_1(\xi_1+1)\}|\leq d^{2\varepsilon}.\]
\item $\alpha_2$ does not contain many jumps:
\[|\{\xi_2\in[2m-1]:\alpha_2(\xi_2)\neq\alpha_2(\xi_2+1)\}|\leq d^{2\varepsilon}.\]
\end{enumerate}
Consider a random walk $W$ on top of the ``chessboard'' and stop it when reaching the boundary of the board (i.e., when it tries to make a step outside the board $[m]\times[2m]$). We define a good step to be a step of the form $(1,1)$ that lands in a black block. We will later relate good steps to violating edges. Our goal is, therefore, to show that a typical $W$ makes many good steps.
\paragraph{Lemma 29.} Let $\delta=0.10$ and assume the chessboard is well-behaved. The probability that $W$ makes less than $d^{2\varepsilon}$ good steps is at most $d^{-\delta}$.
We use this lemma to show Lemma 27.
\paragraph{Proof of Lemma 27 given Lemma 29.} Recall that $A_j$ is an arc of size $|A_j|=m=\lfloor d^{\Delta/(\Delta+1)-3\varepsilon}\rfloor$ so that there is a color $k_j$ satisfying
\[d^{\Delta/(\Delta+1)-4\varepsilon}\leq|S_k\cap A|\leq|A|-d^{\Delta/(\Delta+1)-4\varepsilon}.\tag{2}\]
Furthermore, condition on $P_1,\ldots,P_{\tau_j},|A^{(j-1)}|\geq3$. Assume without loss of generality that $R_{\tau_j}$ is in $A_j$ (when $L_{\tau_j}$ is in $A_j$, the analysis is similar). The distance of $R_{\tau_j}$ from the smallest element of $A_j$ is at most one (the length of ``one step to the right'' is between zero and two). We now grow the random interval until $\sigma_j$, i.e., as long as $R_t$ stays in $A_j$. At the same time, $L_t$ performs a movement to the left. Since $|A^{(j-1)}|\geq3$, there are at least $2m$ steps for $L_t$ to take to the left before hitting $A_j$. There is a one-to-one correspondence between pairings $P$ and random walks $W$ using the correspondence
\begin{align*}
P_{t+1}=\{L_t-2,L_t-1\}&\longleftrightarrow W_{t+1}=W_t+(0,2),\\
P_{t+1}=\{L_t-1,R_t+1\}&\longleftrightarrow W_{t+1}=W_t+(1,1),\\
P_{t+1}=\{R_t+1,R_t+2\}&\longleftrightarrow W_{t+1}=W_t+(2,0).
\end{align*}
Define the function $\alpha_1$ to be $1$ at positions of $A_j$ with color $k_j$, and $0$ at the other positions. Set the function $\alpha_2$ as to describe the color $k_j$ from $L_{\tau_j}$ leftward. The ``chessboard'' is well-behaved by (2) and since $k_j$ is in the set $B$ defined in case 2 of the proof of Lemma 23 (so there are not many jumps for the color $k_j$). Finally, if $W$ makes a good step, then the corresponding pair added to $P$ violated color $k_j$. So, if $E_j$ holds for $P$, then the corresponding $W$ makes less than $d^{2\varepsilon}$ good steps. Formally, by Lemma 23,
\[\mathbb P[E_j\mid P_1,\ldots,P_{\tau_j},|A^{(j-1)}|\geq3]\leq\mathbb P[W\text{ makes less than }d^{2\varepsilon}\text{ good steps}]\leq d^{-\delta}.\]\hfill$\triangleleft$
\paragraph{Proof of Lemma 29.} Define three events $E_R,E_C,E_D$, all of which happen with small probability, so that every $W$ that is not in their union makes many good steps.
Call a subset of the board of the form $I\times[2m]$ or $[m]\times I$, where $I$ is a sub-interval, a region. The width of a region is the size of $I$. Let $R$ be the set of regions of width at least $d^{4\varepsilon}$. The size of $R$ is at most $2m^2$. For a region $r$ in $R$, denote by $E_r$ the event that the number of steps of the form $(1,1)$ that $W$ makes in $r$ is less than $d^{2\varepsilon}$ given that it makes at least $d^{3\varepsilon}$ steps in $r$. Denote
\[E_R=\bigcup_{r\in R}E_r.\]
To estimate the probability of $E_r$, note that we can simply apply the Chernoff bound to a sum of $d^{3\varepsilon}$ Bernoulli random variables with $p=1/3$. By the union bound, we conclude that there is a universal constant $0<c<1$ such that
\[\mathbb P[E_R]\leq c^{d^{3\varepsilon}}.\]
Denote by $H$ the set of all points in the board with $L_1$-norm at least $m^{5/8}$. At time $T$ the random walk $W$ is distributed along all points in $\mathbb N^2$ of $L_1$-norm exactly $T$. The distribution of $W$ on this set is the same as that of a random walk on $\mathbb Z$ that is started at $0$, and moves at every step to the right with probability $1/3$, stays in place with probability $1/3$ and moves to the left with probability $1/3$. The probability that such a random walk on $\mathbb Z$ is at a specific point in $\mathbb Z$ at time $T$ is at most $O(T^{-1/2})$. Hence, for every point $h$ in $H$,
\[\mathbb P[W\text{ hits }h]\leq O(m^{-5/16})\leq m^{-1/4}.\]
Call a point $c=(\xi_1,\xi_2)$ in the board a corner if both $(\xi_1,\xi_2)$ and $(\xi_1+1,\xi_2+1)$ are of the same color $\kappa\in\{\text{black},\text{white}\}$, but $(\xi_1+1,\xi_2)$ and $(\xi_1,\xi_2+1)$ are not of color $\kappa$. For a corner $c$, denote by $\Delta(c)$ the $d^{4\varepsilon}$-neighborhood of $c$ in $L_1$-metric. Denote by $\Delta$ the union over all $\Delta(c)$, for corners $c$ in $H$. Denote by $E_C$ the event that $W$ hits any point in $\Delta$. Since the board is well-behaved, the number of jumps in each of $\alpha_1,\alpha_2$ is at most $d^{2\varepsilon}$. Therefore, the number of corners is at most $d^{4\varepsilon}$. By the union bound,
\[\mathbb P[E_C]\leq O(d^{4\varepsilon}d^{8\varepsilon}m^{-1/4})\leq d^{-0.112},\]
where in the last step, we plugged in $\varepsilon=1/1000$ and used $m\geq d^{1/2-3\varepsilon}$. Next, let $m'=\lceil m^{5/8}\rceil$. Define three (vertical) lines: $D_1$ is the line $\{m'\}\times[2m]$, $D_2$ is the line $\{2m'\}\times[2m]$ and $D_3$ is the line $\{m-m'\}\times[2m]$. Denote by $E_D$ the event that $W$ does not cross the line $D_3$ before stopped (i.e., hitting the boundary of the board). Chernoff's bound implies that there is a universal constant $0<c<1$ for which
\[\mathbb P[E_D]\leq c^m.\]
To conclude the proof by the union bound, it suffices to show that for every $W$ not in $E_R\cup E_C\cup E_D$, the walk $W$ makes at least $d^{2\varepsilon}$ good steps. Fix such a walk $W$. Since $W\notin E_D$, we know that $W$ crosses the line $D_2$.
We consider several cases. Define a block to be a maximal monochromatic rectangle in the board. The board is thus partitioned into black blocks and white blocks -- which is what led [8] to calling it a ``chessboard.'' We now think of the board $[m]\times[2m]$ as drawn in the plane with $(1,1)$ at the bottom-left corner and $(m,2m)$ at the upper-right corner.
\paragraph{Case 1:} The walk $W$ does not hit any white block after crossing $D_1$ and before crossing $D_2$. In this case, all steps taken in the region whose left border is $D_1$ and right border is $D_2$ are in a black area. The number of such steps is at least $m^{5/8}/2\gg d^{3\varepsilon}$. Since $W\notin E_R$, the claim holds.
\paragraph{Case 2:} The walk $W$ hits a white block after crossing $D_1$ and before crossing $D_2$. Let us label the blocks as follows: we associate every block with a pair $\langle\eta_1,\eta_2\rangle$ where $\eta_1$ is between $1$ and the number of jumps in $\alpha_1$ and $\eta_2$ is between $1$ and the number of jumps in $\alpha_2$. So, the label of the ``bottom-left'' is $\langle1,1\rangle$, the label of the block ``above'' it is $\langle1,2\rangle$ and the label of the block ``to its right'' is $\langle2,1\rangle$, etc. There are two sub-cases to consider:
\paragraph{Sub-case 1:} At some point after crossing $D_1$ and before crossing $D_3$, there are two white blocks of the form $\langle\eta_1,\eta_2\rangle$, $\langle\eta_1+1,\eta_2+1\rangle$ so that $W$ intersects both blocks. Let $c$ be the corner between these two blocks (which must exist by definition). Since $W\notin E_C$, we know that $W$ does not visit $\Delta(c)$. Therefore, $W$ must walk in a black area around $\Delta(c)$. Every path surrounding $\Delta(c)$ has length at least $d^{4\varepsilon}$. Since $W\notin E_R$, the claim holds.
\paragraph{Sub-case 2:} At all times after crossing $D_1$ and before crossing $D_3$, the walk never moves from a white block $\langle\eta_1,\eta_2\rangle$ to one of the two white blocks $\langle\eta_1+1,\eta_2+1\rangle$, $\langle\eta_1-1,\eta_2-1\rangle$. Since $W\notin E_D$, this is indeed the last case. The width of a combinatorial rectangle in the board is the size of its ``bottom side'' (i.e., the corresponding subset of $[m]$). Let $\eta$ be the first white block $W$ hits after crossing $D_1$. Let $\Sigma$ be the family of black blocks that are to the right but on the same height as $\eta$. Define $\gamma$ as the maximal width of a rectangle of the form $\sigma\cap[0,m-m_0-1]\times[2m]$ over all $\sigma\in\Sigma$. Since the board is well-behaved, it follows (from the first condition) that the total width of the black area on the same height as $\eta$ is at least $d^{\Delta/(\Delta+1)-4\varepsilon}$. Also, since we are in case 2, the left border of $\eta$ is to the left of $D_2$. Therefore, the total width of the black area to the right of the left border of $\eta$ and to the left of $D_3$, on the same height as $\eta$ is at least $d^{\Delta/(\Delta+1)-4\varepsilon}-3m'$. Therefore, since the number of jumps is at most $d^{2\varepsilon}$,
\[\gamma\geq(d^{\Delta/(\Delta+1)-4\varepsilon}-3m')/d^{2\varepsilon}\gg d^{4\varepsilon}.\]
Since we are in this sub-case, the walk $W$ must ``go through'' every black block it hits: it can go from bottom side to upper side or from left side to right side (but not from left side to upper side or from bottom side to right side). Consider the behaviour of $W$ after it hits $\eta$: starting from a white block, because $W\notin E_D$, it is guaranteed to cross $D_3$. Therefore, the color of the block that $W$ ``exits'' from from each column must keep alternating between white and black. For each black block in $\Sigma$, therefore, there exists a black block in the same column that $W$ crosses horizontally. Focusing on one such black block of width $\gamma$, since $W\notin E_R$, the claim holds.\hfill$\triangleleft$
